\subsection{Local rings}

PROPOSITION 1. Let A be a ring and I the set of non-invertible elements of A. The set I is the union of the ideals of A which are distinct from A. Moreover, the following conditions are equivalent:

\begin{enumerate}
\def\labelenumi{(\alph{enumi})}
\item
  I is an ideal.
\item
  The set of ideals of \(\mathbf{A}\) distinct from \(\mathbf{A}\) has a greatest element.
\item
  A has a unique maximal ideal.
\end{enumerate}

The relation \(x \in I\) is equivalent to \(1 \notin xA\) and hence \(xA \neq A\) . If \(a\) is an ideal of \(A\) distinct from \(A\) and \(x \in a\) , then \(xA \subset a\) , hence \(xA \neq A\) and \(x \in I\) . Hence every ideal of \(A\) distinct from \(A\) is contained in \(I\) and every element \(x \in I\) belongs to a principal ideal \(xA \neq A\) . This proves the first assertion, which immediately implies the equivalence of (a), (b) and (c).

Remark (1). Note that, if (c) holds, I is the Jacobson radical of the ring A (Algebra, Chapter VIII, § 6, no. 3, Definition 3).

DEFINITION 1. A ring A is called a local ring if it satisfies the equivalent conditions (a), (b) and (c) of Proposition 1. The quotient of A by its Jacobson radical (which is then the unique maximal ideal of A) is called the residue field of A.

DEFINITION 2. Let A, B be two local rings and m, n their respective maximal ideals. A homomorphism \(u: \mathbf{A} \to \mathbf{B}\) is called local if \(u(\mathfrak{m}) \subset \mathfrak{n}\) .

This amounts to saying that \(\bar{u}^{1}(\mathfrak{n}) = \mathfrak{m}\) , for \(\bar{u}^{1}(\mathfrak{n})\) is then an ideal containing m and not containing l and hence equal to m. Taking quotients, we then derive canonically from u an injective homomorphism A/m → B/n from the residue field of A to that of B.

\textbf{Examples}\par

\begin{enumerate}
\def\labelenumi{(\arabic{enumi})}
\item
  A field is a local ring. A ring which is reduced to 0 is not a local ring.
\item
  Let A be a local ring and k its residue field. The ring of formal power series \(B = A[[X_{1}, \ldots, X_{n}]]\) is a local ring, for the non-invertible elements of B are the formal power series whose constant terms are not invertible in A (Algebra, Chapter IV, § 5, no. 6, Proposition 4). The canonical injection of A into B is a local homomorphism and the corresponding injection of residue fields is an isomorphism.
\item
  Let \(\mathfrak{b}\) be an ideal of a ring A which is only contained in a single maximal ideal m; then A/b is a local ring with maximal ideal m/b and residue field canonically isomorphic to A/m. This applies in particular to the case \(b = m^{k}\) , m being any maximal ideal of A ( \(\S\) 1, no. 1, Corollary to Proposition 1). If A itself is a local ring with maximal ideal m, then for every ideal \(b \neq A\) of A, A/b is a local ring, the canonical homomorphism \(A \to A/b\) a local homomorphism and the corresponding homomorphism of residue fields an isomorphism.
\item
  Let X be a topological space, \(x_{0}\) a point of X and A the ring of germs at the point \(x_{0}\) of real-valued functions continuous in a neighbourhood of \(x_{0}\) (General Topology, Chapter I, §6, no. 10). Clearly, for the germ at \(x_{0}\) of a continuous function f to be invertible in A, it is necessary and sufficient that \(f(x_{0}) \neq 0\) , as that implies that \(f(x) \neq 0\) in a neighbourhood of \(x_{0}\) . The ring A is therefore a local ring whose maximal ideal m is the set of germs of functions which are zero at \(x_{0}\) ; taking quotients, the mapping \(g \mapsto g(x_{0})\) of A to R gives an isomorphism of the residue field A/m onto R.
\end{enumerate}

PROPOSITION 2. Let A be a ring and p a prime ideal of A. The ring \(A_{p}\) is local; its maximal ideal is the ideal \(pA_{p} = p_{p}\) , generated by the canonical image of p in \(A_{p}\) ; its residue field is canonically isomorphic to the field of fractions of A/p.

Let \(S = A - p\) and \(j: A \to A_p\) be the canonical homomorphism; the hypothesis that \(p\) is prime implies that \(p\) is saturated with respect to \(S\) , hence \(j^{1}(\mathfrak{pA}_{\mathfrak{p}}) = \mathfrak{p}\) (§ 2, no. 4, Proposition 10) and, as the ideals of A not meeting S are those contained in \(\mathfrak{p}\) , the first two assertions are special cases of § 2, no. 5, Proposition 11 (ii). Moreover, if \(f\) is the canonical homomorphism \(A \to A/\mathfrak{p}\) , \(f(S)\) is the set of elements \(\neq 0\) of the integral domain \(A/\mathfrak{p}\) and hence the last assertion is a special case of § 2, no. 5, Proposition 11 (i).

DEFINITION 3. Let A be a ring and p a prime ideal of A. The ring \(A_{p}\) is called the local ring of A at p, or the local ring of p, when there is no ambiguity.

Remark (2). If A is a local ring and m its maximal ideal, the elements of A --- m are invertible (Proposition 1) and hence \(A_{m}\) is canonically identified with A (§ 2, no. 1, Remark 5).

\textbf{Examples}\par

\begin{enumerate}
\def\labelenumi{(\arabic{enumi})}
\setcounter{enumi}{4}
\tightlist
\item
  Let \(p\) be a prime number. The local ring \(\mathbf{Z}_{(p)}\) is the set of rational numbers \(a / b\) , where \(a, b\) are rational integers with \(b\) prime to \(p\) ; the residue field of \(\mathbf{Z}_{(p)}\) is isomorphic to the prime field \(\mathbf{F}_p = \mathbf{Z} / (p)\) .
\end{enumerate}

* (6) Let V be an affine algebraic variety, A the ring of regular functions on V, W an irreducible subvariety of V and p the (necessarily prime) ideal of A consisting of the functions which are zero at every point of W. The ring \(\mathbf{A}_{\mathfrak{p}}\) is called the local ring of W on V. *

PROPOSITION 3. Let A be a ring, p a prime ideal of A and S = A - p.~For every ideal \(\mathfrak{b}'\) of \(\mathbf{A}_{\mathfrak{p}}\) distinct from \(\mathbf{A}_{\mathfrak{p}}\) , let \(\mathfrak{b}\) be the ideal \((i_{\mathbf{A}}^{\mathbf{S}})^{-1}(\mathfrak{b}')\) of A so that \(\mathfrak{b}' = \mathfrak{bA}_{\mathfrak{p}}\) .

\begin{enumerate}
\def\labelenumi{(\roman{enumi})}
\item
  Let \(f\) be the canonical homomorphism \(\mathbf{A} \to \mathbf{A} / \mathfrak{b}\) . The homomorphism from \(\mathbf{A}_{\mathfrak{p}}\) to \((\mathbf{A} / \mathfrak{b})_{\mathfrak{p} / \mathfrak{b}}\) canonically associated with \(f\) (§ 2, no. 1, Proposition 2) is surjective and its kernel is \(\mathfrak{b}'\) , which defines, by taking quotients, a canonical isomorphism of \(\mathbf{A}_{\mathfrak{p}} / \mathfrak{b}'\) onto \((\mathbf{A} / \mathfrak{b})_{\mathfrak{p} / \mathfrak{b}}\) .
\item
  The mapping \(\mathfrak{b}' \to \mathfrak{b} = (i_{\mathbf{A}}^{\mathbb{S}})^{-1}(\mathfrak{b}')\) , restricted to the set of prime ideals of \(\mathbf{A}_{\mathfrak{p}}\) , is an isomorphism (with respect to inclusion) of this set onto the set of prime ideals of \(\mathbf{A}\) contained in \(\mathfrak{p}\) . If \(\mathfrak{b}'\) is prime in \(\mathbf{A}_{\mathfrak{p}}\) , there exists an isomorphism of the ring \(\mathbf{A}_{\mathfrak{b}}\) onto the ring \((\mathbf{A}_{\mathfrak{p}})_{\mathfrak{b}'}\) , which maps \(a / s\) to \((a / 1) / (s / 1)\) for all \(a \in \mathbf{A}\) , \(s \in \mathbf{A} - \mathfrak{b}\) ).
\end{enumerate}

This is just a particular case of § 2, no. 5, Proposition 11.

\textbf{Remarks}\par

\begin{enumerate}
\def\labelenumi{(\arabic{enumi})}
\setcounter{enumi}{2}
\item
  If \(\mathfrak{a}\) is an ideal of \(\mathbf{A}\) not contained in \(\mathfrak{p}\) , then \(\mathfrak{a}\mathrm{A}_{\mathfrak{p}} = \mathrm{A}_{\mathfrak{p}}\) and \((\mathrm{A} / \mathfrak{a})_{\mathfrak{p}} = 0\) (§ 2, no. 5, Remark).
\item
  Let A, B be two rings, \(\rho: A \to B\) a homomorphism, \(q\) a prime ideal of B and \(p\) the prime ideal \(\rho^{-1}(q)\) of A. As \(\rho(A - p) \subset B - q\) , a canonical homomorphism \(\rho_q: A_p \to B_q\) is derived from \(\rho\) ( \(\S 2\) , no. 1, Proposition 2) and it is immediate that \(\rho_q(pA_p) \subset qB_q\) , hence \(\rho_q\) is a local homomorphism.
\end{enumerate}
% semantic-fragments: legacy-input-seam
\relax{}
